% Laget automatisk av elevark/lag_elevark.py fra opplegg/intro-bevegelsesmengde.html. Endringer her blir overskrevet.
\documentclass[11pt]{article}
\usepackage{elevark}
\begin{document}
\elevarkhode{Fysikk 1 · Bevegelsesmengde og støt}{Bevegelsesmengde og impuls – arbeidsark}{https://fysikklab.no/sim/impuls.html}{Myk landing}

\section*{Oppgaver}

Alle oppgavene bruker Myk landing\footnote{\url{fysikklab.no/sim/impuls.html}}. Velg positiv retning oppover. Newtons 2. lov kan skrives som $ΣF \cdot Δt = Δp$, og impulsen $ΣF \cdot Δt$ er arealet under grafen for $F$ mot $t$.

\begin{oppgaveliste}
\oppgave \fase{gjett}{Gjett}Egget slippes fra 1,0~m. Blir endringen i bevegelsesmengde større på betong enn på en dyne? Blir kraften større?
\svar
\oppgave \fase{test}{Test}Slipp egget fra 1,0~m på alle de fem underlagene. Skriv opp bremsetiden, $Δp$ og den største kraften i en tabell. Hva er likt, og hva er forskjellig?
\svar
\oppgave \fase{forklar}{Forklar}Velg «Felles for alle landinger» som skala, og se på grafene for betong og dyne. Hva kan du si om arealet under grafene fram til egget står stille? Bruk $ΣF \cdot Δt = Δp$ til å forklare hvorfor egget knuser på betong.
\svar
\oppgave \fase{test}{Regn}Regn ut farten og bevegelsesmengden rett før egget lander fra 1,0~m. Egget lander på skumgummi og stopper på 17~ms. Regn ut middelkraften fra underlaget. Husk at tyngden også virker.
\svar
\oppgave \fase{gjett}{Gjett}Velg sprettball og betong. Ballen har samme masse (60~g) og samme fart før landingen som egget. Blir $Δp$ større, mindre eller like stor som for egget? Begrunn svaret før du slipper.
\svar
\oppgave \fase{test}{Regn}Slipp ballen fra 1,0~m. Les av farten den spretter opp med, og regn ut $p$ før, $p$ etter og $Δp$. Sjekk svaret med pilene og tallene i simuleringen.
\svar
\oppgave \fase{forklar}{Forklar}Hvorfor blir $Δp$ nesten dobbelt så stor for ballen som for egget? Slipp ballen på alle underlagene. Mellom hvilke to verdier ligger $Δp$ alltid?
\svar
\end{oppgaveliste}

\end{document}
